\IfFileExists{stacks-project.cls}{%
\documentclass{stacks-project}
}{%
\documentclass{amsart}
}

% For dealing with references we use the comment environment
\usepackage{verbatim}
\newenvironment{reference}{\comment}{\endcomment}
%\newenvironment{reference}{}{}
\newenvironment{slogan}{\comment}{\endcomment}
\newenvironment{history}{\comment}{\endcomment}

% For commutative diagrams we use Xy-pic
\usepackage[all]{xy}

% We use 2cell for 2-commutative diagrams.
\xyoption{2cell}
\UseAllTwocells

% We use multicol for the list of chapters between chapters
\usepackage{multicol}

% This is generally recommended for better output
\usepackage{lmodern}
\usepackage[T1]{fontenc}

% For cross-file-references
\usepackage{xr-hyper}

% Package for hypertext links:
\usepackage{hyperref}

% For any local file, say "hello.tex" you want to link to please
% use \externaldocument[hello-]{hello}
\externaldocument[introduction-]{introduction}
\externaldocument[conventions-]{conventions}
\externaldocument[sets-]{sets}
\externaldocument[categories-]{categories}
\externaldocument[topology-]{topology}
\externaldocument[sheaves-]{sheaves}
\externaldocument[sites-]{sites}
\externaldocument[stacks-]{stacks}
\externaldocument[fields-]{fields}
\externaldocument[algebra-]{algebra}
\externaldocument[brauer-]{brauer}
\externaldocument[homology-]{homology}
\externaldocument[derived-]{derived}
\externaldocument[simplicial-]{simplicial}
\externaldocument[more-algebra-]{more-algebra}
\externaldocument[smoothing-]{smoothing}
\externaldocument[modules-]{modules}
\externaldocument[sites-modules-]{sites-modules}
\externaldocument[injectives-]{injectives}
\externaldocument[cohomology-]{cohomology}
\externaldocument[sites-cohomology-]{sites-cohomology}
\externaldocument[dga-]{dga}
\externaldocument[dpa-]{dpa}
\externaldocument[sdga-]{sdga}
\externaldocument[hypercovering-]{hypercovering}
\externaldocument[schemes-]{schemes}
\externaldocument[constructions-]{constructions}
\externaldocument[properties-]{properties}
\externaldocument[morphisms-]{morphisms}
\externaldocument[coherent-]{coherent}
\externaldocument[divisors-]{divisors}
\externaldocument[limits-]{limits}
\externaldocument[varieties-]{varieties}
\externaldocument[topologies-]{topologies}
\externaldocument[descent-]{descent}
\externaldocument[perfect-]{perfect}
\externaldocument[more-morphisms-]{more-morphisms}
\externaldocument[flat-]{flat}
\externaldocument[groupoids-]{groupoids}
\externaldocument[more-groupoids-]{more-groupoids}
\externaldocument[etale-]{etale}
\externaldocument[chow-]{chow}
\externaldocument[intersection-]{intersection}
\externaldocument[pic-]{pic}
\externaldocument[weil-]{weil}
\externaldocument[adequate-]{adequate}
\externaldocument[dualizing-]{dualizing}
\externaldocument[duality-]{duality}
\externaldocument[discriminant-]{discriminant}
\externaldocument[derham-]{derham}
\externaldocument[local-cohomology-]{local-cohomology}
\externaldocument[algebraization-]{algebraization}
\externaldocument[curves-]{curves}
\externaldocument[resolve-]{resolve}
\externaldocument[models-]{models}
\externaldocument[functors-]{functors}
\externaldocument[equiv-]{equiv}
\externaldocument[pione-]{pione}
\externaldocument[etale-cohomology-]{etale-cohomology}
\externaldocument[proetale-]{proetale}
\externaldocument[relative-cycles-]{relative-cycles}
\externaldocument[more-etale-]{more-etale}
\externaldocument[trace-]{trace}
\externaldocument[crystalline-]{crystalline}
\externaldocument[spaces-]{spaces}
\externaldocument[spaces-properties-]{spaces-properties}
\externaldocument[spaces-morphisms-]{spaces-morphisms}
\externaldocument[decent-spaces-]{decent-spaces}
\externaldocument[spaces-cohomology-]{spaces-cohomology}
\externaldocument[spaces-limits-]{spaces-limits}
\externaldocument[spaces-divisors-]{spaces-divisors}
\externaldocument[spaces-over-fields-]{spaces-over-fields}
\externaldocument[spaces-topologies-]{spaces-topologies}
\externaldocument[spaces-descent-]{spaces-descent}
\externaldocument[spaces-perfect-]{spaces-perfect}
\externaldocument[spaces-more-morphisms-]{spaces-more-morphisms}
\externaldocument[spaces-flat-]{spaces-flat}
\externaldocument[spaces-groupoids-]{spaces-groupoids}
\externaldocument[spaces-more-groupoids-]{spaces-more-groupoids}
\externaldocument[bootstrap-]{bootstrap}
\externaldocument[spaces-pushouts-]{spaces-pushouts}
\externaldocument[spaces-chow-]{spaces-chow}
\externaldocument[groupoids-quotients-]{groupoids-quotients}
\externaldocument[spaces-more-cohomology-]{spaces-more-cohomology}
\externaldocument[spaces-simplicial-]{spaces-simplicial}
\externaldocument[spaces-duality-]{spaces-duality}
\externaldocument[formal-spaces-]{formal-spaces}
\externaldocument[restricted-]{restricted}
\externaldocument[spaces-resolve-]{spaces-resolve}
\externaldocument[formal-defos-]{formal-defos}
\externaldocument[defos-]{defos}
\externaldocument[cotangent-]{cotangent}
\externaldocument[examples-defos-]{examples-defos}
\externaldocument[algebraic-]{algebraic}
\externaldocument[examples-stacks-]{examples-stacks}
\externaldocument[stacks-sheaves-]{stacks-sheaves}
\externaldocument[criteria-]{criteria}
\externaldocument[artin-]{artin}
\externaldocument[quot-]{quot}
\externaldocument[stacks-properties-]{stacks-properties}
\externaldocument[stacks-morphisms-]{stacks-morphisms}
\externaldocument[stacks-limits-]{stacks-limits}
\externaldocument[stacks-cohomology-]{stacks-cohomology}
\externaldocument[stacks-perfect-]{stacks-perfect}
\externaldocument[stacks-introduction-]{stacks-introduction}
\externaldocument[stacks-more-morphisms-]{stacks-more-morphisms}
\externaldocument[stacks-geometry-]{stacks-geometry}
\externaldocument[moduli-]{moduli}
\externaldocument[moduli-curves-]{moduli-curves}
\externaldocument[examples-]{examples}
\externaldocument[exercises-]{exercises}
\externaldocument[guide-]{guide}
\externaldocument[desirables-]{desirables}
\externaldocument[coding-]{coding}
\externaldocument[obsolete-]{obsolete}
\externaldocument[fdl-]{fdl}
\externaldocument[index-]{index}

% Theorem environments.
%
\theoremstyle{plain}
\newtheorem{theorem}[subsection]{Theorem}
\newtheorem{proposition}[subsection]{Proposition}
\newtheorem{lemma}[subsection]{Lemma}

\theoremstyle{definition}
\newtheorem{definition}[subsection]{Definition}
\newtheorem{example}[subsection]{Example}
\newtheorem{exercise}[subsection]{Exercise}
\newtheorem{situation}[subsection]{Situation}

\theoremstyle{remark}
\newtheorem{remark}[subsection]{Remark}
\newtheorem{remarks}[subsection]{Remarks}

\numberwithin{equation}{subsection}

% Macros
%
\def\lim{\mathop{\mathrm{lim}}\nolimits}
\def\colim{\mathop{\mathrm{colim}}\nolimits}
\def\Spec{\mathop{\mathrm{Spec}}}
\def\Hom{\mathop{\mathrm{Hom}}\nolimits}
\def\Ext{\mathop{\mathrm{Ext}}\nolimits}
\def\SheafHom{\mathop{\mathcal{H}\!\mathit{om}}\nolimits}
\def\SheafExt{\mathop{\mathcal{E}\!\mathit{xt}}\nolimits}
\def\Sch{\mathit{Sch}}
\def\Mor{\mathop{\mathrm{Mor}}\nolimits}
\def\Ob{\mathop{\mathrm{Ob}}\nolimits}
\def\Sh{\mathop{\mathit{Sh}}\nolimits}
\def\NL{\mathop{N\!L}\nolimits}
\def\CH{\mathop{\mathrm{CH}}\nolimits}
\def\proetale{{pro\text{-}\acute{e}tale}}
\def\etale{{\acute{e}tale}}
\def\QCoh{\mathit{QCoh}}
\def\Ker{\mathop{\mathrm{Ker}}}
\def\Im{\mathop{\mathrm{Im}}}
\def\Coker{\mathop{\mathrm{Coker}}}
\def\Coim{\mathop{\mathrm{Coim}}}

% Boxtimes
%
\DeclareMathSymbol{\boxtimes}{\mathbin}{AMSa}{"02}

%
% Macros for moduli stacks/spaces
%
\def\QCohstack{\mathcal{QC}\!\mathit{oh}}
\def\Cohstack{\mathcal{C}\!\mathit{oh}}
\def\Spacesstack{\mathcal{S}\!\mathit{paces}}
\def\Quotfunctor{\mathrm{Quot}}
\def\Hilbfunctor{\mathrm{Hilb}}
\def\Curvesstack{\mathcal{C}\!\mathit{urves}}
\def\Polarizedstack{\mathcal{P}\!\mathit{olarized}}
\def\Complexesstack{\mathcal{C}\!\mathit{omplexes}}
% \Pic is the operator that assigns to X its picard group, usage \Pic(X)
% \Picardstack_{X/B} denotes the Picard stack of X over B
% \Picardfunctor_{X/B} denotes the Picard functor of X over B
\def\Pic{\mathop{\mathrm{Pic}}\nolimits}
\def\Picardstack{\mathcal{P}\!\mathit{ic}}
\def\Picardfunctor{\mathrm{Pic}}
\def\Deformationcategory{\mathcal{D}\!\mathit{ef}}
\usepackage{cleveref}
\crefname{originaldefinition}{original definition}{original definitions}
\Crefname{originaldefinition}{Original definition}{Original definitions}
\makeatletter
\newlabel{situation-complexes}{{0DLE}{source}{}{sourcepointer-1}{}}
\newlabel{situation-complexes@cref}{{[originaldefinition][0][]0DLE}{[1][0][]source}{}{}{}}
\makeatother
\crefname{originallemma}{original lemma}{original lemmas}
\Crefname{originallemma}{Original lemma}{Original lemmas}
\makeatletter
\newlabel{algebraic-lemma-representable-diagonal}{{045G}{source}{}{sourcepointer-2}{}}
\newlabel{algebraic-lemma-representable-diagonal@cref}{{[originallemma][0][]045G}{[1][0][]source}{}{}{}}
\makeatother
\crefname{originallemma}{original lemma}{original lemmas}
\Crefname{originallemma}{Original lemma}{Original lemmas}
\makeatletter
\newlabel{lemma-complexes}{{0DLD}{source}{}{sourcepointer-3}{}}
\newlabel{lemma-complexes@cref}{{[originallemma][0][]0DLD}{[1][0][]source}{}{}{}}
\makeatother
\crefname{originallemma}{original lemma}{original lemmas}
\Crefname{originallemma}{Original lemma}{Original lemmas}
\makeatletter
\newlabel{lemma-complexes-open-neg-exts-vanishing}{{0DLC}{source}{}{sourcepointer-4}{}}
\newlabel{lemma-complexes-open-neg-exts-vanishing@cref}{{[originallemma][0][]0DLC}{[1][0][]source}{}{}{}}
\makeatother
\crefname{originalproposition}{original proposition}{original propositions}
\Crefname{originalproposition}{Original proposition}{Original propositions}
\makeatletter
\newlabel{proposition-isom}{{08K9}{source}{}{sourcepointer-5}{}}
\newlabel{proposition-isom@cref}{{[originalproposition][0][]08K9}{[1][0][]source}{}{}{}}
\makeatother
\title{The diagonal of the complexes stack is representable by algebraic spaces}
\author{The Stacks Project Authors}
\date{stacks-project:quot.tex:local-snapshot-20261012; GNU FDL 1.2 or later}
\begin{document}
\maketitle
\setcounter{subsection}{1}
\noindent\textit{Editorial note.} All delivered mathematical text is the author's own, copied from the pinned Stacks Project chapter file quot.tex; no statement or proof step was rewritten, reordered or removed. The counted span is the author environment and its own proof environment. References inside the span are handed to the assembler as navigation stubs that carry the canonical tag of the cited tag only; whether each of them is load-bearing is recorded on the stub itself.\par\medskip
\noindent\textit{Allowed standard prerequisites (editorial).}\begin{itemize}
\item Situation situation-complexes (canonical tag 0DLE). cited at statement line 5427. Standing situation.
\item Lemma algebraic-lemma-representable-diagonal (canonical tag 045G). cited at proof lines 5436-5437. Goal reformulation: Isom over T is an algebraic space.
\item Lemma lemma-complexes (canonical tag 0DLD). cited at proof lines 5450-5451. E, E' satisfy conditions (1), (2).
\item Lemma lemma-complexes-open-neg-exts-vanishing (canonical tag 0DLC). cited at proof lines 5462-5463 and 5476-5477. Opens W, W' and affine fp representability of the SheafHom functors.
\item Proposition proposition-isom (canonical tag 08K9). cited at proof lines 5482-5483. The Isom-as-fibre-product presentation.
\end{itemize}\medskip
\ifdefined\hypertarget\hypertarget{sourcepointer-1}{}\fi
\noindent\textit{Source reference.} Original definition 0DLE: sources/stacks-chow/tags.txt, source lines 17044--17044. Cited at statement line 5427: the lemma is phrased in this situation; load-bearing.\par\medskip
\ifdefined\hypertarget\hypertarget{sourcepointer-2}{}\fi
\noindent\textit{Source reference.} Original lemma 045G: sources/stacks-chow/tags.txt, source lines 5435--5435. Cited at proof lines 5436-5437: the goal reformulation; load-bearing.\par\medskip
\ifdefined\hypertarget\hypertarget{sourcepointer-3}{}\fi
\noindent\textit{Source reference.} Original lemma 0DLD: sources/stacks-chow/tags.txt, source lines 17043--17043. Cited at proof lines 5450-5451: E, E' satisfy conditions (1), (2); load-bearing.\par\medskip
\ifdefined\hypertarget\hypertarget{sourcepointer-4}{}\fi
\noindent\textit{Source reference.} Original lemma 0DLC: sources/stacks-chow/tags.txt, source lines 17042--17042. Cited at proof lines 5462-5463 and 5476-5477 (twice): the opens W, W' and the representability of H = SheafHom(E, E') and friends; load-bearing twice.\par\medskip
\ifdefined\hypertarget\hypertarget{sourcepointer-5}{}\fi
\noindent\textit{Source reference.} Original proposition 08K9: sources/stacks-chow/tags.txt, source lines 10877--10877. Cited at proof lines 5482-5483: repeating its argument gives Isom(E, E') = (H' x\_B H) x\_\{c, I x\_B I', sigma\} B; load-bearing.\par\medskip
\begin{lemma}
\label{lemma-complexes-diagonal}
In Situation \ref{situation-complexes}. Denote
$\mathcal{X} = \Complexesstack_{X/B}$. Then
$\Delta : \mathcal{X} \to \mathcal{X} \times \mathcal{X}$ is
representable by algebraic spaces.
\end{lemma}

\begin{proof}
Consider two objects $x = (T, g, E)$ and $y = (T, g', E')$
of $\mathcal{X}$ over a scheme $T$. We have to show that
$\mathit{Isom}_\mathcal{X}(x, y)$ is an algebraic space over $T$, see
Algebraic Stacks, Lemma \ref{algebraic-lemma-representable-diagonal}.
If for $h : T' \to T$ the restrictions $x|_{T'}$ and $y|_{T'}$ are isomorphic
in the fibre category $\mathcal{X}_{T'}$, then $g \circ h = g' \circ h$.
Hence there is a transformation of presheaves
$$
\mathit{Isom}_\mathcal{X}(x, y) \longrightarrow \text{Equalizer}(g, g')
$$
Since the diagonal of $B$ is representable (by schemes) this equalizer is
a scheme. Thus we may replace $T$ by this equalizer and
$E$ and $E'$ by their pullbacks. Thus we may assume $g = g'$.

\medskip\noindent
Assume $g = g'$. After replacing $B$ by $T$ and $X$ by $X_T$ we arrive
at the following problem. Given $E, E' \in D(\mathcal{O}_X)$
satisfying conditions (1), (2) of Lemma \ref{lemma-complexes}
we have to show that $\mathit{Isom}(E, E')$ is an algebraic space.
Here $\mathit{Isom}(E, E')$ is the functor
$$
(\Sch/B)^{opp} \to \textit{Sets},\quad
T \mapsto \{\varphi : E_T \to E'_T
\text{ isomorphism in }D(\mathcal{O}_{X_T})\}
$$
where $E_T$ and $E'_T$ are the derived pullbacks of $E$ and $E'$
to $X_T$. Now, let $W \subset B$, resp.\ $W' \subset B$ be the
open subspace of $B$ associated
to $E, E'$, resp.\ to $E', E$ by
Lemma \ref{lemma-complexes-open-neg-exts-vanishing}.
Clearly, if there exists an isomorphism $E_T \to E'_T$ as in
the definition of $\mathit{Isom}(E, E')$, then we see that
$T \to B$ factors into both $W$ and $W'$ (because we have
condition (1) for $E$ and $E'$ and we'll obviously have
$E_t \cong E'_t$ so no nonzero maps $E_t[i] \to E_t$
or $E'_t[i] \to E_t$ over the fibre $X_t$ for $i > 0$.
Thus we may replace $B$ by the open $W \cap W'$.
In this case the functor $H = \SheafHom(E, E')$
$$
(\Sch/B)^{opp} \to \textit{Sets},\quad T \mapsto
\Hom_{\mathcal{O}_{X_T}}(E_T, E'_T)
$$
is an algebraic space affine and of finite presentation over $B$ by
Lemma \ref{lemma-complexes-open-neg-exts-vanishing}.
The same is true for
$H' = \SheafHom(E', E)$,
$I = \SheafHom(E, E)$, and
$I' = \SheafHom(E', E')$.
Therefore we can repeat the argument of the proof of
Proposition \ref{proposition-isom} to see that
$$
\mathit{Isom}(E, E') = (H' \times_B H) \times_{c, I \times_B I', \sigma} B
$$
for some morphisms $c$ and $\sigma$. Thus
$\mathit{Isom}(E, E')$ is an algebraic space.
\end{proof}

\end{document}
